\title{LongRoPE: Extending LLM Context Window Beyond 2 Million Tokens}

\begin{abstract}
Expanding the context window is a sought-after capability for large language models (LLMs). Nevertheless, advancements are hampered by the substantial cost of fine-tuning, the limited availability of extended textual data, and the disruptive effects caused by the introduction of new token positions. As a result, existing large context windows generally do not exceed approximately 128k tokens.

This work presents LongRoPE, a method that, for the first time, expands the context window of pre-trained LLMs to a remarkable \textbf{2048k} tokens. This is accomplished with no more than 1k fine-tuning steps using training sequences up to 256k tokens in length, all while preserving the model’s performance on its original short context window. Our approach is based on three main contributions: (i) through an efficient search process, we discover and leverage two kinds of non-uniformities within positional interpolation, improving fine-tuning initialization and supporting up to an 8$\times$ extension in contexts without fine-tuning; (ii) we propose a staged extension procedure, whereby an LLM is initially fine-tuned on 256k token sequences, after which an additional round of positional interpolation is performed on the resulting extended LLM, enabling a 2048k token window; (iii) LongRoPE is recalibrated at an 8k context length to restore its performance on short contexts. Comprehensive experiments on LLaMA2 and Mistral models spanning multiple tasks validate the efficacy of the proposed methodology. LongRoPE-augmented models maintain their original architectures, requiring only slight adjustments to the positional embeddings, and are compatible with most existing optimizations. Code will be released at \texttt{https://github.com/microsoft/LongRoPE}

\section{Introduction}

Large Language Models (LLMs) have achieved impressive results across a range of tasks~\cite{gpt4,llama2}, but they are still constrained by a restricted context window—for instance, LLaMA2 is limited to 4096 tokens~\cite{llama2}. When sequences exceed this predefined window, model performance deteriorates since LLMs are not exposed to these extended positions during training. This limitation hinders applications such as in-context learning involving many examples~\cite{Huang2023FewerIM} as well as LLM-based agent systems~\cite{park2023generative,madaan2023selfrefine}.

Recent studies have demonstrated that it is possible to extend the context window of a pre-trained LLM up to approximately 128k tokens through fine-tuning with longer input texts~\cite{longlora,pi,yarn,zhang2024soaring,liu2023scaling}. Nonetheless, advancing beyond this limit presents three significant challenges. Firstly, new position indices, when not previously seen during training, can generate erratic values that result in out-of-distribution complications, thus impeding the fine-tuning process and hindering convergence~\cite{pi}. This difficulty is amplified during drastic expansions, for instance, from 4k to over $>$1000k positions, where over 90\% of position indices are novel. Secondly, fine-tuning typically depends on the availability of sufficiently long training samples. Current datasets offer a scarcity of ultra-long documents, particularly those surpassing 1000k tokens. Further, training with these extended sequences imposes substantial computational costs, necessitating extensive processing time and large-scale GPU allocation. Thirdly, scaling the context window to extremely high token counts diffuses the attention mechanism, dispersing it over a vast number of positions. This dilution leads to a notable decline in model performance when handling the shorter contexts seen in the original pre-training phase~\cite{pi}.

A commonly adopted strategy for addressing the first challenge involves interpolating the RoPE positional embeddings~\cite{rope,pi}, whereby position indices beyond those seen during pretraining are rescaled to fall within the original range (see Fig.~\ref{fig:overview}). The Position Interpolation (PI) method~\cite{pi} performs a linear interpolation of RoPE's rotational angles using an extension ratio. In contrast, NTK~\cite{ntk,dynamicntk} employs non-uniform interpolation and extrapolation techniques across different RoPE dimensions. YaRN~\cite{yarn} further refines this process by sorting RoPE dimensions into three groups according to their frequencies, subsequently applying extrapolation, NTK-based, and linear interpolation schemes to each group. Despite these advancements, positional embeddings in Transformer models are characterized by a \emph{complex and spatially-varying information entropy}, a nuance inadequately addressed by current interpolation techniques. As a result, these methods fail to capitalize on this heterogeneity, which in turn contributes to information degradation and constrains the possible length of the context window.

Section~\ref{sec:analysis} presents two principal empirical insights: \textbf{(1)} Successful positional interpolation needs to address two types of non-uniformities, namely, the diversity among RoPE dimensions and differences in token positions. Reduced interpolation is advantageous for lower RoPE dimensions and for tokens at earlier positions, but the precise configuration should be tailored to the intended extrapolation length. \textbf{(2)} When these non-uniformities are incorporated into positional interpolation, the core information encoded by the original RoPE—especially in essential dimensions and positions—is better preserved. This approach mitigates interpolation-induced information loss, resulting in improved starting points for subsequent fine-tuning. Additionally, it enables up to an 8$\times$ increase in context length even without further fine-tuning.

Building on these insights, we present \textbf{LongRoPE}, a robust approach designed to push the LLM context window to more than 2 \emph{million} tokens. LongRoPE hinges on three principal advancements. The first is that {LongRoPE} takes full advantage of multidimensional non-uniformity within positional interpolation. Specifically, it determines suitable rescale factors for the rotation angles in RoPE for each dimension, conditioned on the token positions. Given that the space for identifying optimal rescale factors grows exponentially with the desired extension ratio, {LongRoPE} employs an evolutionary search algorithm, enhanced by two specialized optimization strategies, to improve the efficiency of the search. Figure~\ref{fig:overview} presents an illustration of the resulting rescaled RoPE discovered through this method.

Subsequently, {LongRoPE} implements a resource-efficient, gradually expanding method to attain a 2048k context window, circumventing the necessity for direct fine-tuning on ultra-long input sequences, which are rare and difficult to procure. Initially, the approach involves identifying an optimal 256k context length on the pre-trained LLM, followed by fine-tuning at this length. Next, exploiting the capability of our non-uniform positional interpolation to facilitate up to an 8$\times$ context extension in settings that do not require additional fine-tuning, we perform a further search to determine updated RoPE rescale factors using the previously fine-tuned, length-extended LLM. This sequence of steps enables {LongRoPE} to ultimately reach a 2048k context window for both LLaMA2 and Mistral~\cite{mistral}.

In order to address the issue of reduced performance on the standard (short) context window, {LongRoPE} maintains adjustments to the RoPE rescale factors within the enlarged LLM. Just as the model transitions from a 256k to a 2048k window, we also downscale to 4k and 8k context lengths on the LLM fine-tuned at 256k, applying our search algorithm to discourage extensive positional interpolation. At inference time, for any sequence shorter than 8k, RoPE is modified to use the rescale factors obtained through this search procedure.

Comprehensive experimentation involving multiple LLMs and a range of long-context tasks substantiates the efficacy of our approach. The results indicate that LongRoPE consistently preserves low perplexity across evaluation lengths ranging from 4k up to 2048k, attains more than 90\% accuracy in passkey retrieval, and achieves accuracy on par with existing methods in standard benchmarks constructed within a 4096 context window. LongRoPE is broadly compatible with all LLMs utilizing RoPE embeddings. We plan to make our code and the LongRoPE-2048k models publicly available.

\section{Non-uniformity in Positional Interpolation}
\label{sec:analysis}

\subsection{Preliminary}

Transformer architectures depend on clear positional cues, typically supplied as position embeddings, to capture the sequence of input tokens. In this study, we concentrate on RoPE~\cite{rope} position embeddings, a prevalent choice in state-of-the-art LLMs. The RoPE encoding for a token located at position $n$ can be succinctly expressed as:

\begin{equation}
		\fontsize{7.8}{7.8} \selectfont
	\label{eq:rope}
	\left[ cos(n\theta_0), sin(n\theta_0), cos(n\theta_1), \cdot\cdot\cdot, cos(n\theta_{d/2-1}), sin(n\theta_{d/2-1}) \right]   
\end{equation}

Here, $d$ denotes the dimension of the embedding, 
$n\theta_i$ corresponds to the rotary angle associated with the token in position $n$, and 
$\theta_i = \theta^{-2i/d}$ captures the rotation frequencies used. In the context of RoPE, the parameter $\theta$ is set to a default value of $10000$.

\textbf{\emph{Context window extension ratio $s$} and positional interpolation}. We define $s$ as  the ratio of extended context length $L'$ to the original length $L$: $s=\frac{L'}{L}$. 

To increase the context window from $L$ to $L'$, existing positional interpolation approaches propose adjusting the rotation frequencies $\theta_i$ downward according to the expansion ratio $s$. Introducing $\beta=\theta^{2/d}$ and letting $\lambda$ represent the precise scaling factor connected to $s$, we present a unified formulation for these positional interpolation strategies as:

\begin{equation}	
	\fontsize{7.5}{7.5} \selectfont
	\label{eq:unified_rope}
	\begin{aligned}
		\left[\cos\left(\frac{n}{\lambda(\beta)^0}\right), \sin \left(\frac{n}{\lambda(\beta)^0}\right), \cos\left(\frac{n}{\lambda(\beta)^1}\right), \ldots,  \sin \left(\frac{n}{\lambda(\beta)^{d/2-1}}\right) \right]   
	\end{aligned}
\end{equation}

\textbf{Linear positional interpolation (PI)}. PI~\cite{pi} proposes a linear rescaling of position indices within the range established during pre-training. When adapting to an extended sequence length determined by a scaling factor $s$, the method linearly decreases the rotation angles for each position by a factor of $\lambda=s$ in every RoPE dimension. Nonetheless, this compression of position information leads to densely packed positional encodings, which impairs the model's capacity to differentiate among nearby tokens. As a result, PI often yields suboptimal results when the extension ratio becomes large.

\textbf{NTK-based interpolation and extrapolation}. ~\cite{ntk,dynamicntk} examine RoPE through the lens of information encoding, employing Neural Tangent Kernel (NTK) theory~\cite{jacot2018neural,tancik2020fourier} for their analysis. To address the issue of positional overcrowding seen in PI, these works propose distributing the interpolation burden across the various RoPE dimensions. Specifically, they adjust scaling such that lower (higher frequency) dimensions are scaled down less, while higher (lower frequency) dimensions are scaled up more. This leads to both improved positional interpolation and enables extrapolation, parameterized by $\lambda=s^{i}$. The enhanced dynamic NTK~\cite{dynamicntk} further adapts the extension ratio at each sequence position according to the actual sequence length. In contrast to PI, which requires fine-tuning, approaches informed by NTK are capable of context window extension in the absence of fine-tuning, though they are generally limited to a maximum extension ratio of 4$\times$.

\textbf{YaRN}~\cite{yarn} offers a notable enhancement in the effectiveness of positional interpolation. The method segments the RoPE dimensions into three groups based on frequency, applying distinct interpolation approaches to each. Specifically, extrapolation ($\lambda$=1) is assigned to high frequency dimensions, while linear interpolation (PI) is reserved for those with low frequencies. Intermediate frequencies are addressed using the NTK scheme. A central aspect of YaRN is this partitioning of RoPE dimensions; however, the determination of these groups currently relies on heuristic, human-driven procedures. Consequently, this reliance may hinder optimal results when adapting to novel LLM architectures.

\subsection{Study on Non-uniform Positional Interpolation}

Drawing on insights from NTK and YaRN, it becomes apparent that their improvements originate from introducing non-linearity—particularly by leveraging distinct frequencies along RoPE dimensions to better facilitate targeted interpolation and extrapolation. Yet, existing methods for non-linearity predominantly depend on hand-crafted heuristics. This observation prompts two central questions: (1) Is the prevailing approach to positional interpolation truly optimal? (2) Might there exist alternative, untapped non-linear functions?

In order to address these questions, we employ evolution search (as described in Sec.~\ref{sec:method}) to identify improved non-uniform positional interpolations for LLaMA2-7B. This search process leverages perplexity as a metric, evaluating with 5 randomly selected samples from the PG19~\cite{pg19} validation dataset. Our experimental investigation yields several important observations.

\textbf{Finding 1}: \textit{RoPE dimensions exhibit substantial non-uniformities, which are not effectively handled by current positional interpolation methods.}

We conduct a search for the optimal $\lambda$ value corresponding to each RoPE dimension as specified in Eq.~\ref{eq:unified_rope}. Table~\ref{tbl:exp1} presents a perplexity comparison of LLaMA2-7B evaluated with various approaches on the PG19 and Proof-pile~\cite{proof-pile} test sets, with all models evaluated in a zero-shot manner (i.e., without additional fine-tuning). Our learned configuration achieves substantial gains over existing techniques, indicating that prevailing strategies—both linear (PI) and non-uniform (Dynamic-NTK and YaRN) interpolations—do not fully optimize performance. Interestingly, we observe that YaRN actually lags behind PI and NTK on the PG19 benchmark. This can be attributed to YaRN failing to consistently cover the intended context window length when applied to language models that have not undergone further fine-tuning. For instance, when using an 8k context window, perplexity with YaRN rises sharply beyond the 7k token mark.

Our exploration reveals that the rescaling coefficients $\lambda$ in Eq.~\ref{eq:unified_rope} are no longer uniform, unlike the constant scale $s$ employed by PI, the computation of NTK, and the group-level scaling of YaRN. These variable scaling factors lead to notable improvements in LLaMA2's language modeling capabilities—specifically in terms of perplexity—when using context windows of 8k and 16k tokens, without necessitating additional fine-tuning. This improvement arises because the adapted positional embedding maintains much of the structure of the original RoPE, with an emphasis on preserving crucial dimensions. As a result, the LLM faces fewer challenges when differentiating between closely positioned tokens.

\textbf{Finding 2}: \textit{RoPE for the initial tokens in the input sequence should be extrapolated with less interpolation.} 

We propose that RoPE should reduce interpolation for the initial $\hat{n}$ tokens in the input, as these tokens generally draw higher attention scores and are thus particularly important for the attention mechanism—consistent with findings from Streaming LLM~\cite{streamingllm} and LM-Infinite~\cite{han2023lminfinite}. To test this idea, we conduct experiments extending the context length to 8k and 16k via PI and NTK, while systematically excluding the first $\hat{n}$ tokens (with $\hat n$ taking values 0, 2, ..., 256) from interpolation. Note that setting $\hat n=0$ corresponds to using the standard versions of PI and NTK without modification. The results, summarized in Table~\ref{tbl:tokenposition}, yield two primary insights: \textbf{(1)} skipping position interpolation for the sequence's earliest tokens leads to noticeable improvements in model performance; and \textbf{(2)} the ideal value for $\hat{n}$ varies depending on the desired context extension length.

\textbf{Finding 3}: \textit{Non-uniform positional interpolation effectively extends LLM context window in both fine-tuning and non-fine-tuning settings.} 

Although our implementation of non-uniform position interpolation demonstrates notable gains in extension capabilities at 8k and 16k context lengths without further training, achieving successful performance on longer contexts necessitates fine-tuning. Therefore, we apply fine-tuning to LLaMA2-7B using our optimized RoPE configuration for a 64k context window (refer to the Appendix for details on the fine-tuning procedure). As illustrated in Table~\ref{tbl:finetune}, our approach yields a marked performance advantage over both PI and YaRN, both in the pre-fine-tuned and post-fine-tuned states of LLaMA2-7B. This improvement can be attributed to our method's strategic adoption of non-uniform positional interpolation, which reduces the extent of information loss and ensures more robust initialization for the fine-tuning phase.

\textbf{Summary}. Through our analysis, we identify two sources of non-uniformity—differences in RoPE dimensions and token positions. Leveraging these non-uniform factors within positional interpolation enables substantial enhancements in large language model context window extension.

\section{LongRoPE}

\label{sec:method}
Building on these insights, we propose LongRoPE—a method that initially develops an efficient search strategy to leverage both forms of non-uniformity, and subsequently applies this approach to scale the LLM context window past 2 million tokens.

\subsection{Problem Formulation}

The presence of these two sources of non-uniformity results in an expansive solution space and introduces significant challenges for optimization procedures. To mitigate these issues, we recast the multidimensional non-uniform position interpolation task as a search-based optimization problem.

Consider an LLM configured for a context window of size $L'$, handling large input documents $\mathbf{X}$ with each sequence $\mathbf{x} \in \mathbf{X}$ containing more tokens than $L'$. In this context, we define the original rotary angle in the $i^{th}$ RoPE embedding dimension at token position $n$ by $\frac{n}{\beta^i}$. Accordingly, we can state the optimization objective as:

\begin{equation}
\begin{gathered}
		\label{eq:problem}

\underset {\mathbf{x}\in\mathbf{X}; \, |\mathbf{x}|\ge L'}{ \operatorname {arg\,min} } \,  \mathcal{L}\, \left( \text{LLM} (\text{RoPE}, \mathbf{X})\right), \text{where}
\\
\scalebox{0.93}{
$\underset {\begin{subarray}{l} i=0,\ldots,\frac{d}{2}-1;\\ n\in[0,|\mathbf{x}|);\end{subarray}}{\text{RoPE}_{(n)}} = {\left[\,\cdots,\, \cos\Big( \mathbb{I}(\hat{\lambda}_i, \hat{n}) \frac{n}{\beta^i} \Big),\, \sin\Big( \mathbb{I}(\hat{\lambda}_i, \hat{n}) \frac{n}{\beta^i} \Big),\,\cdots \,\right]}$}\\
	\text{with } \mathbb{I}(\hat{\lambda}_i, \hat{n}) = \begin{cases}
	1&\text{if $n < \hat{n}$} \\
	\frac{1}{\lambda_i}&\text{if $n \geq \hat{n}$}
	\end{cases}
	\end{gathered}
\end{equation}

We incorporate a collection of rescale factors, denoted as $\mathbb{I}(\hat{\lambda}_{i},\hat{n})$, to address the two types of non-uniformity considered. Here, $\hat{\lambda}_i$ captures variations across RoPE dimensions, while $\hat{n}$ represents non-uniformities along token positions. In particular, $\mathbb{I}(\hat{\lambda}_{i},\hat{n})$ is applied to adjust the rotation angle for the $i^{th}$ dimension within RoPE; the factor $\hat{\lambda}_{i}$ serves as a scaling coefficient, and $\hat{n}$ establishes the token position threshold beyond which rescaling occurs. For the first $\hat{n}$−1 token positions, the scaling is inactive and the standard RoPE rotational angle $\frac{n}{\beta^i}$ remains unchanged. When token positions satisfy $n\ge\hat{n}$, the introduced rescale factor becomes active and modifies the rotary angle accordingly.

With a desired context window size of $L'$, our goal is to determine the set of optimal rescale factors ($\mathbb{I}(\hat{\lambda}_{0},\hat{n})$, $\mathbb{I}(\hat{\lambda}_{1},\hat{n})$, ...$\mathbb{I}(\hat{\lambda}_{i},\hat{n})$...) across the $1^{st}$ through $d^{th}$ dimensions of RoPE. By applying these rescaling coefficients to the target $\text{LLM}$’s $\text{RoPE}$, we seek to minimize the next token prediction loss $\mathcal{L}$ (that is, perplexity) when processing input samples $\mathbf{X}$ whose sequence length is $L'$.

\subsection{Searching the Non-uniform Position Interpolation}

To address the challenge posed in Eq.~\ref{eq:problem}, we present a straightforward yet powerful approach that identifies the most suitable RoPE rescale factors, enabling us to take full advantage of the multidimensional non-uniformity present in position embeddings.

\textbf{Search space}. Our method defines an expansive search space that incorporates both forms of non-uniformity. This space is detailed in Table~\ref{tbl:searchspace}. In particular, we permit independent optimization of a rescale factor for each dimension within the RoPE embeddings. For greater simplicity in search space formulation, we choose to search over $\lambda_i$ and $\hat{n}$, rather than directly searching for $\mathbb{I}(\hat{\lambda}_{i},\hat{n})$, where  $\hat{\lambda}_{i}=1/\lambda_i$. 
As outlined in Table~\ref{tbl:searchspace}, the search interval for $\lambda_i$ ranges from a minimum of 1.0 (corresponding to pure extrapolation) up to a maximum of $s\times1.25$ (which enables interpolation beyond standard PI), with increments of 0.01. Here, $s$ denotes the intended proportion for extending the context window.

The parameter $\hat{n}$ determines how many of the initial token positions are preserved with their original RoPE embedding, excluding them from position interpolation. In practice, we consider $\hat{n}$ values within the set \{0, 1, 2, 4, 8, 12, 16, 20, 24, 28, 32, 64, 128, 256\} during the search process. If $\hat{n}=0$, the searched rescale factors are applied to all token positions.

\textbf{Evolution-based search}. The search space described in Table~\ref{tbl:searchspace} encompasses a vast array of positional interpolation configurations, making efficient navigation through it particularly demanding. As an illustration, selecting an extension rate of $s=4\times$ results in $400^{128/2}\times14=4\times10^{167}$ possible combinations. Notably, increasing the extension ratio further magnifies the search space exponentially. To tackle this complexity, we adopt evolution search~\cite{guo2020single} and propose two optimization strategies that considerably enhance the efficiency of the search process. The complete search methodology is outlined in Algorithm~\ref{alg:search}.

\textit{Optimized initial population generation}. Rather than fully randomizing all $P$ rescale factors in the initial population, we explicitly include the three RoPE rescale factors associated with PI, NTK, and YaRN as predefined individuals. The other $P$-3 individuals are created by applying random mutations to these three rescale factors, each with a probability of $p$.

\textit{Monotonically non-decreasing constraint}. Once the initial population has been established, LLM perplexity values are calculated for all individuals. To do so, the respective RoPE rescale factors are implemented within the target LLM, and perplexity is then measured for the input $\mathbf{X}$. The highest-performing $k$ individuals are subsequently selected as parent candidates for the next evolutionary step. Nevertheless, due to the enormity of the search space, undirected application of mutation and crossover operators may result in suboptimal candidate solutions and excess perplexity evaluations. This inefficiency becomes even more pronounced when $L'$ is large, as each perplexity measurement is computationally expensive.

To tackle this issue, a monotonic non-decreasing constraint is enforced on the set of sampled RoPE rescaled factors: $\lambda_i \leq \lambda_{i+1}$. Only those RoPE configurations that adhere to this condition are utilized when evaluating perplexity in LLMs, which leads to a notable reduction in search complexity. In particular, we stipulate that the values of $\lambda_i$ must increase monotonically with respect to the RoPE dimension (where $i=0,...,63$). The motivation for imposing this dimensional monotonicity draws from NTK theory~\cite{jacot2018neural,tancik2020fourier,ntk}, which indicates that lower dimensions—associated with higher frequencies—necessitate less interpolation (smaller $\lambda_i$), while higher dimensions—corresponding to lower frequencies—are capable of supporting more interpolation (larger $\lambda_i$).

\begin{algorithm}[t]
	\small
	\caption{The search algorithm}
	\textbf{Input:} target LLM, input samples $\mathbf{X}$, population size $P$, mutation size $N_1$, crossover size $N_2$, max iterations $\mathcal{T}$, mutation probability $p$\\

	\begin{algorithmic}[1]
		\label{alg:search}
		\STATE Initialize Top-k as the empty set;
		\STATE $\text{P}_0$ = \textit{Initialize\_population\_with\_optimization} ($P$, $\mathbf{X}$, $p$);
		\FOR{$i=1$ \textbf{to} $\mathcal{T}$}
		\STATE \textit{Compute\_perplexity} (LLM, $\text{P}_{i-1}$, $\mathbf{X}$);
		\STATE Top-k $\leftarrow$ \textit{Update\_Topk} (Top-k, $\text{P}_{i-1}$);
		\STATE $\text{P}_{mutation}$ = \textit{Mutation\_with\_mono\_constraint}(Top-k, $N_1$, $p$);
		\STATE $\text{P}_{crossover}$ = \textit{Crossover\_with\_mono\_constraint}(Top-k, $N_2$);
		\STATE $\text{P}_i$ = $\text{P}_{mutation}$ $\cup$ $\text{P}_{crossover}$ $\cup$ Top-k;
		\ENDFOR
		\STATE Output the member in Top-k achieving minimal perplexity;
	\end{algorithmic}
\end{algorithm}

\textbf{8$\times$ extension without fine-tuning}. Through evolutionary search, we successfully discover optimal non-uniform RoPE rescaling factors that retain critical dimensions and position information, thereby reducing information degradation due to interpolation. As shown in Fig.\ref{fig:non-ft}, our approach enables LLaMA2's context window to scale from 4k up to 32k tokens without the need for fine-tuning. Comparatively, other techniques—including PI, as well as non-uniform NTK and YaRN—experience significant increases in perplexity after merely doubling the context length.

\subsection{Extending LLM Context Window to 2048K}
\label{sec:extendto2048k}

\textbf{Progressive extension to 2048k}. In this section, we present our approach for expanding the context window of pre-trained LLMs from the conventional 4k up to beyond 2048k. Our results indicate that non-uniform positional interpolation permits an 8$\times$ increase in context length without any requirement for fine-tuning. However, when greater extensions such as 512$\times$ are needed, fine-tuning becomes indispensable. A common strategy involves identifying suitable RoPE rescaling parameters for the intended 2048k window followed by fine-tuning. Nevertheless, this process is hampered by the immense computational and training costs involved. Additionally, our empirical observations reveal that effective fine-tuning for such substantial extension ratios is highly challenging (see Appendix).

Fortunately, LongRoPE demonstrates efficacy with both the base model and the fine-tuned variant extended for larger context windows. As such, we propose a resource-efficient, incremental approach that attains the target length of 2048k in only 1k fine-tuning steps, confining the training context length to a maximum of 256k.

$\Diamond$ \textit{LongRoPE search applied to expanding pre-trained LLM to 256k}. Using LLaMA2 for illustration, we perform a search corresponding to desired context window lengths of 128k and 256k. At this phase, the extension factors are 32$\times$ and 64$\times$, accordingly.

$\Diamond$ \textit{Fine-tuning to 256k}. Next, we adapt the pre-trained LLM to support a context window of 256k tokens through fine-tuning. Our approach involves initially fine-tuning LLaMA2 for 400 steps using RoPE rescaled factors corresponding to 128k. After this stage, we update the finished checkpoint by applying the 256k RoPE rescaled factors, and continue with a further 600 fine-tuning steps. This staged procedure offers greater efficiency compared to directly fine-tuning with 256k settings from the outset.

$\Diamond$ \textit{Scaling fine-tuned extended LLM up to 2048k with LongRoPE search}. In the final step, we apply a subsequent search procedure to the already fine-tuned LLM at 256k sequence length. This enables us to expand the context window up to 2048k tokens, all without the need for any additional fine-tuning. The achieved overall extension factor is $512\times$.

\textbf{Restoring performance in shorter context windows}. When scaling up the context window to a substantial 2048k length, we observe diminished performance within the language model's initial context span. This degradation is a recognized drawback of positional interpolation~\cite{pi}, which compacts the position embeddings for the original context window into a limited area of the higher-dimensional space, impairing model effectiveness. Notably, with an extension factor of 512$\times$, the embeddings corresponding to the first 4k positions are densely packed, exacerbating this issue.

To address this issue, we conduct an additional evolutionary search on the extended LLM to fine-tune the RoPE rescale factors specifically for shorter context lengths such as 4k and 8k tokens. For these reduced lengths, we lower the upper limit of the searched $\lambda$ parameter, since less positional interpolation is needed. When running inference, the LLM automatically modifies the relevant RoPE rescale factors in response to context length.

\section{LongRoPE}

\label{sec:method}
Building on these insights, we propose LongRoPE, which initially incorporates an efficient search algorithm designed to leverage both identified non-uniformities. Leveraging this approach, we successfully expand the LLM context window to surpass 2 million tokens.

\subsection{Problem Formulation}

The presence of these two forms of non-uniformity significantly expands the solution space and increases the difficulty of the optimization process. To overcome this challenge, we recast the multidimensional non-uniform position interpolation optimization challenge as a search problem.

Consider an LLM designed for a context window of size $L'$ and a set of lengthy documents $\mathbf{X}$, where each individual sequence $\mathbf{x} \in \mathbf{X}$ contains more tokens than $L'$. We represent the standard rotary angle for the $i^{th}$ dimension within the RoPE embedding at the $n^{th}$ token position as $\frac{n}{\beta^i}$. Based on these definitions, the optimization task can be formally specified as follows:

\begin{equation}
\begin{gathered}
		\label{eq:problem}

\underset {\mathbf{x}\in\mathbf{X}; \, |\mathbf{x}|\ge L'}{ \operatorname {arg\,min} } \,  \mathcal{L}\, \left( \text{LLM} (\text{RoPE}, \mathbf{X})\right),\quad \text{with}
\\
\scalebox{0.93}{
$\underset {\begin{subarray}{l} i=0,\cdot\cdot,\frac{d}{2}-1;\\ n\in[ 0,|\mathbf{x}|);\end{subarray}}{\text{RoPE}_(n)}= {\left[\ldots, \cos\left(\mathbb{I}(\hat{\lambda}_{i}, \hat{n})\cdot\frac{n}{\beta^i}\right),  \sin\left(\mathbb{I}(\hat{\lambda}_{i}, \hat{n})\cdot\frac{n}{\beta^i}\right),\ldots\right] }$}
\\
\text{where}\quad \mathbb{I}(\hat{\lambda}_{i},\hat{n})=\begin{cases}
	1&\mbox{if $n<\hat{n}$}\\
{\frac{1}{\lambda}_{i}}&\mbox{if $n\geq\hat{n}$}
\end{cases}
	\end{gathered}
\end{equation}

In this framework, we define a group of rescaling factors, $\mathbb{I}(\hat{\lambda}_{i},\hat{n})$, designed to address the two distinct types of non-uniformity. Here, $\hat{\lambda}_i$ characterizes the non-uniformity across RoPE dimensions, while $\hat{n}$ pertains to non-uniformity related to token positions. Concretely, $\mathbb{I}(\hat{\lambda}_{i},\hat{n})$ serves to adjust the rotation angle in the $i^{th}$ RoPE dimension, with $\hat{\lambda}_i$ functioning as the scaling parameter and $\hat{n}$ as the specified positional threshold. For the first $\hat{n} - 1$ tokens, the rescaling via $\hat{\lambda}_i$ is omitted, so the standard RoPE rotary angle $\frac{n}{\beta^i}$ is retained. Only when token positions satisfy $n \geq \hat{n}$ does the rescaling factor become operative.

Assuming a desired context window length of $L'$, our task is to identify the set of optimal rescale factors ($\mathbb{I}(\hat{\lambda}_{0},\hat{n})$, $\mathbb{I}(\hat{\lambda}_{1},\hat{n})$, ... , $\mathbb{I}(\hat{\lambda}_{i},\hat{n})$, ...) across each RoPE dimension from $1^{st}$ to $d^{th}$. By applying these rescale factors to the $\text{RoPE}$ in the target $\text{LLM}$, our aim is to minimize the next token prediction loss, denoted as $\mathcal{L}$ (corresponding to perplexity), over input sequences $\mathbf{X}$ of token length $L'$.

\subsection{Searching the Non-uniform Position Interpolation}

In order to address the challenge described in Eq.~\ref{eq:problem}, we propose a straightforward but powerful approach that identifies the most suitable RoPE rescale factors to maximize the benefits from multidimensional non-uniformities present in position embeddings.

\textbf{Search space}. We construct an extensive search space that accommodates both types of non-uniformity. The design of this space is outlined in Table~\ref{tbl:searchspace}. Notably, our method permits the determination of a unique rescale factor for every dimension within the RoPE embeddings. For simplicity in search space construction, we opt to search over $\lambda_i$ and $\hat{n}$, as opposed to directly optimizing for $\mathbb{I}(\hat{\lambda}_{i},\hat{n})$, where $\hat{\lambda}_{i}$ is defined as $1/\lambda_i$. According to Table~\ref{tbl:searchspace}, the allowed range for $\lambda_i$ extends from a lower bound of 1.0 (equivalent to naive extrapolation) to an upper bound of $s\times1.25$ (enabling broader interpolation compared to PI), incrementing in steps of 0.01. Here, $s$ represents the ratio by which the context window is to be expanded.

The parameter $\hat{n}$ determines how many of the first token positions maintain their original configuration, foregoing position interpolation and instead utilizing the standard RoPE embedding. In practice, we consider $\hat{n}$ values from the set \{0, 1, 2, 4, 8, 12, 16, 20, 24, 28, 32, 64, 128, 256\}. Setting $\hat{n}=0$ indicates that position interpolation with the optimized rescale factors is applied to every token position.

\textbf{Evolution-based search}. The search space detailed in Table~\ref{tbl:searchspace} encompasses a vast array of positional interpolation configurations, resulting in substantial complexity for effective exploration. For instance, an extension ratio of $s=4\times$ yields a total of $400^{128/2}\times14=4\times10^{167}$ possible options. As the extension ratio increases, the size of the search space rises exponentially. To manage this complexity, we employ evolution search~\cite{guo2020single} and propose two additional optimization methods to significantly enhance search efficiency. The complete search framework is outlined in Algorithm~\ref{alg:search}.

\textit{Optimized initial population generation}. Rather than relying solely on random initialization for the population of $P$ rescale factors, our approach incorporates the three RoPE rescale factors associated with PI, NTK, and YaRN directly into the starting population as individual members. The rest of the $P$-3 individuals are produced by randomly applying mutations to these three rescale factors with probability $p$.

\textit{Monotonically non-decreasing constraint}. Following the initialization of the population, we evaluate the LLM perplexity for each member. This involves applying the specific RoPE rescale factors to the chosen LLM and determining the perplexity of the input $\mathbf{X}$. The top-k candidates are selected as parents for the subsequent evolutionary steps. Nonetheless, the extensive search space means that simple mutation and crossover operations can frequently yield low-quality solutions, resulting in redundant perplexity evaluations. This inefficiency becomes more pronounced as $L'$ increases, since computing each perplexity value is computationally intensive.

To tackle this issue, we enforce a constraint of non-decreasing monotonicity on the selected RoPE rescale factors: $\lambda_i \le \lambda_{i+1}$. Only those RoPE configurations adhering to this restriction are considered in the LLM perplexity assessment, thereby markedly reducing the computational demands of the search process. In more detail, the imposed condition requires the sequence of $\lambda_i$ values to be monotonically increasing across the RoPE dimensions (with $i=0,\ldots, 63$). The rationale for enforcing dimensional monotonicity originates from NTK theory~\cite{jacot2018neural,tancik2020fourier,ntk}, which indicates that dimensions associated with higher frequencies (lower indices) necessitate less interpolation (corresponding to a smaller $\lambda_i$), whereas dimensions capturing lower frequencies (higher indices) are able to accommodate greater levels of interpolation (a larger $\lambda_i$).

\begin{algorithm}[t]
	\small
	\caption{The search algorithm}
	\textbf{Input:} target LLM, input samples $\mathbf{X}$, population size $P$, mutation size $N_1$, crossover size $N_2$, maximum iterations $\mathcal{T}$, mutation probability $p$\\

	\begin{algorithmic}[1]
		\label{alg:search}
		\STATE Top-k $\gets$ $\phi$;
		\STATE $\text{P}_0 \gets$ \textit{Initialize\_population\_with\_optimization}($P$, $\mathbf{X}$, $p$);
		\FOR{$i = 1$ \TO $\mathcal{T}$}
		\STATE \textit{Compute\_perplexity}(LLM, $\text{P}_{i-1}$, $\mathbf{X}$);
		\STATE Top-k $\gets$ \textit{Update\_Topk}(Top-k, $\text{P}_{i-1}$);
		\STATE $\text{P}_{mutation} \gets$ \textit{Mutation\_with\_mono\_constraint}(Top-k, $N_1$, $p$);
		\STATE $\text{P}_{crossover} \gets$ \textit{Crossover\_with\_mono\_constraint}(Top-k, $N_2$);
		\STATE $\text{P}_i \gets \text{P}_{mutation} \cup \text{P}_{crossover} \cup$ Top-k;
		\ENDFOR
		\STATE Return the member in Top-k that yields the minimum perplexity;
	\end{algorithmic}
\end{algorithm}

\textbf{8$\times$ extension without fine-tuning}. By leveraging evolutionary search, our approach efficiently discovers optimal, non-uniform RoPE rescale parameters that selectively preserve crucial dimensions and token positions, thus reducing the loss of information typically associated with interpolation. As shown in Fig.\ref{fig:non-ft}, this enables us to expand LLaMA2's context length from 4k up to 32k tokens with no need for further model fine-tuning. By comparison, current techniques such as PI, as well as non-uniform NTK and YaRN, experience a significant rise in perplexity after extending the context window only by 2$\times$.

\subsection{Extending LLM Context Window to 2048K}
\label{sec:extendto2048k}

\textbf{Progressive extension to 2048k}. In this section, we present our approach for expanding the context window of existing pre-trained LLMs from the standard 4k up to more than 2048k tokens. As previously shown, the proposed non-uniform positional interpolation allows for an 8$\times$ increase without necessitating fine-tuning. When even greater expansions are targeted (such as 512$\times$), fine-tuning becomes essential. A common strategy involves identifying suitable RoPE rescale factors tailored for the desired 2048k context and then performing fine-tuning. Nevertheless, this approach is hindered by the extremely high computational demands of training at such scales. Furthermore, our observations indicate that effectively fine-tuning LLMs with a very large extension ratio is nontrivial and presents significant difficulties (see Appendix for further discussion).

Fortunately, LongRoPE demonstrates efficacy with both the base model and the extended fine-tuned LLM. As such, we propose a resource-efficient, incremental approach that reaches the desired 2048k context window using only 1k fine-tuning iterations and a maximum training sequence length of 256k.

$\Diamond$ \textit{LongRoPE search for adapting pre-trained LLM to 256k}. Using LLaMA2 as a case study, we perform a search procedure targeting context window lengths of 128k and 256k tokens. At this phase, the context window is expanded by a factor of 32$\times$ and 64$\times$, respectively.

$\Diamond$ \textit{Fine-tuning to 256k}. Subsequently, the pre-trained LLM is fine-tuned in order to support a context window of 256k. Concretely, LLaMA2 is initially fine-tuned for 400 steps utilizing the RoPE rescaled factors designed for 128k. Afterwards, these factors are substituted with those for 256k on the resulting checkpoint, followed by a further 600 steps of fine-tuning. Compared to direct fine-tuning for 256k, this staged approach demonstrates greater efficiency.

$\Diamond$ \textit{Scaling fine-tuned extended LLM to 2048k via LongRoPE search}. In the final step, we apply a subsequent search method to the LLM previously fine-tuned on sequences of 256k tokens. Through this procedure, we are able to achieve an exceptionally large context window of 2048k, eliminating the need for any additional fine-tuning. The context window thus expands by a factor of $512\times$.

\textbf{Recovery for shorter context windows}. Upon scaling up to a substantial 2048k context window, we observe diminished performance inside the initial context window. This phenomenon is well-documented in studies of positional interpolation~\cite{pi}, where position embeddings for the original context window are compressed into a smaller subspace due to higher-dimensional scaling. Such compression adversely impacts the language model's capabilities. Specifically, with an extension ratio of 512$\times$, the representation of positions in the initial 4k context window becomes heavily congested.

To address this issue, we conduct an additional evolution search on the expanded LLM to fine-tune the RoPE rescale factors specifically for shorter context lengths such as 4k and 8k. For these shorter lengths, the search space for the maximum $\lambda$ is constrained since less positional interpolation is necessary. During inference, the LLM adapts the relevant RoPE rescale factors in real-time.

 \section{Experiments}

\label{sec:exp}
\subsection{Setup}
\textbf{Evaluation Tasks and Models}. LongRoPE is integrated with both LLaMA2-7B and Mistral-7B models. To assess performance, we consider three dimensions: (1) perplexity on lengthy texts for language models operating with extended context; (2) a Passkey retrieval challenge designed to test the model's capacity to locate a specific passkey amid substantial distractor content; and (3) conventional LLM benchmark tasks restricted to a context window of 4096 tokens.

\textbf{Fine-tuning}. In the case of LLaMA2, we set the learning rate at $2\times 10^{-5}$ employing a linear decay schedule, alongside a total batch size of 32. The model undergoes 400 steps of fine-tuning on the Redpajama~\cite{redpajama} dataset, with input sequences divided into 128k-token chunks, each surrounded by BOS and EOS tokens. Subsequently, utilizing the resulting checkpoint, a further 600 steps are conducted to extend the context window to 256k. The training with a 128k context window is carried out on 8 A100 GPUs utilizing the distributed training system~\cite{cube}, whereas the expanded 256k context requires 16 A100 GPUs. For Mistral, the setup uses a fixed learning rate of $1\times 10^{-6}$, a total batch size of 64, and, for both the 128k and 256k variants, adopts the YaRN~\cite{yarn} training configuration. Specifically, we run 400 steps on the Together Computer's Long-Data Collections~\cite{mistraldata}, using sequences of 16k length, and utilize 4 A100 GPUs throughout training.

\textbf{Search}. For target window sizes up to 256k, the configuration is as follows: $P$=64, $N_1$=$N_2$=16, $p$=0.3, $\mathcal{T}$=40, selecting the top 32 candidates for mutation and crossover at each step. Perplexity evaluation utilizes 5 randomly chosen validation samples from the PG19 dataset, each meeting a minimum context length equal to the specified target. For window sizes exceeding 512k, we reduce the population, mutation, and crossover counts by half. In these cases, perplexity is assessed on 3 randomly drawn samples from the Pile-Books3~\cite{pile} validation set.

\textbf{Baselines}. In extending the context window to 2048k, we fine-tune models at intermediate window sizes of 128k and 256k, resulting in LongRoPE-2048k (ft=128k) for LLaMA2 and LongRoPE-2048k (ft=256k) for Mistral. Our analysis benchmarks these four models against leading methods for expanding context length, focusing on open-source LLMs that underwent fine-tuning subsequent to positional interpolation using PI, NTK, and YaRN. The comparative models encompass Together-32k~\cite{together}, Code LLaMA~\cite{codellama}, LongLoRA-full-FT-100k~\cite{longlora}, as well as YaRN-LLaMA and YaRN-Mistral~\cite{yarn}.

\subsection{Main Results}

\label{sec:mainresult}
\textbf{Language modeling on long sequences up to 256k}. To start, we benchmark our approach against leading extended LLMs on sequences as long as 256k tokens. Our experiments utilize two datasets to assess the robustness of our method: the Proof-pile~\cite{pg19} and PG19~\cite{pile} test partitions. We calculate perplexity across a range of context lengths, employing a sliding window of 256 tokens. For PG19, evaluation is conducted on its full test set, comprising 100 documents. In the case of Proof-pile, following the protocol established by YaRN~\cite{yarn}, we sample 10 examples at random, each containing at least 128k tokens.

Table~\ref{tbl:proof-pile} and Table~\ref{tbl:pg19} present a comparison of perplexity scores for LLaMA2 and Mistral models that have been augmented using various interpolation strategies, evaluated on Proof-pile and PG19. Two main findings emerge from these results: \textbf{(1)} the augmented models exhibit a consistent decline in perplexity as the evaluation length increases from 4k to 256k, which demonstrates their capacity to effectively utilize extended context lengths. \textbf{(2)} Despite operating within a context window that is 16$\times$ larger than the original, a condition known to impair short-context performance, our LongRoPE-2048k models still surpass leading baselines when evaluated on contexts of up to 256k tokens.

\textbf{Language modeling on sequences exceeding 2000k tokens}. To assess performance on exceptionally lengthy texts, we utilize the Books3~\cite{pile} dataset. For practical evaluation purposes, we randomly sample 20 books with lengths greater than 2048k tokens, applying a sliding window approach of 256k.

According to the results presented in Table~\ref{tbl:books3}, LongRoPE effectively increases the context window of both LLaMA2-7B and Mistral-7B models to 2048k, while maintaining perplexity at levels that match or surpass baseline methods for shorter input lengths (8k–128k). Noteworthy differences arise between 2048k versions of LLaMA2 and Mistral models: Mistral achieves superior performance compared to baselines at lower context lengths, yet its perplexity rises above 7 when processing beyond 256k tokens. For LLaMA2, as anticipated, perplexity tends to decrease with the ability to leverage longer contexts, aside from modest rises observed at 1024k and 2048k lengths. Additionally, on LLaMA2, LongRoPE-2048k demonstrates better outcomes when fine-tuned on 256k-token windows than with 128k, attributable to a reduced secondary extension ratio (namely, 8$\times$ compared to 16$\times$). In contrast, Mistral obtains higher performance at a 128k fine-tuning window. This difference can be traced to Mistral’s training configuration—specifically, using a training window of 16k for 128k and 256k fine-tuning tasks, in accordance with YaRN, which in turn impacts the model’s capacity to further expand its context window after fine-tuning.

\textbf{Passkey retrieval}. Next, we examine the practical context window limitations within generation tasks. We adopt the synthetic passkey retrieval evaluation introduced by~\cite{passkey}, where a model is tasked with identifying a randomly generated passkey (a five-digit numeral) embedded within an extensive document. The specific prompt format is described in the appendix. For evaluation, we execute the passkey retrieval procedure 10 times, each time positioning the passkey at a random spot uniformly chosen within the total length of the context used for assessment.

Fig.~\ref{fig:passkey} presents a comparison of retrieval accuracy against several baselines. The accuracy of previous models declines precipitously to zero beyond 128k tokens. Conversely, LongRoPE-LLaMA2-2048k (ft=256k) sustains notably high retrieval accuracy ($\ge$90\%) across a broad span from 4k up to 2048k, even in the extremely difficult scenario of extracting a passkey from a context containing millions of tokens. Similarly, LongRoPE-Mistral-2048k (ft=128k) achieves perfect (100\%) retrieval accuracy through 1800k tokens and experiences a drop to 60\% at 2048k, which corresponds to the trend observed in Table~\ref{tbl:books3} where the perplexity rises modestly at 2048k.

\textbf{Evaluation on standard benchmarks within the native context length}. LongRoPE-2048k models are assessed on their performance within the original context window utilizing the Hugging Face Open LLM Leaderboard~\cite{huggingfaceboard}, across both zero-shot and few-shot paradigms. Specifically, the evaluation involves 25-shot ARC-Challenge~\cite{allenai:arc}, 10-shot HellaSwag~\cite{zellers2019hellaswag}, 5-shot MMLU~\cite{mmlu}, and 0-shot TruthfulQA~\cite{lin2021truthfulqa}.

As indicated in Table~\ref{tbl:commonsense}, our models perform similarly to baseline scores on the initial benchmark, which was developed for a shorter context window, and even surpass the original Mistral on TruthfulQA by a margin of +0.5\%. Although LongRoPE-LLaMA2-2048k, which underwent fine-tuning at 256k, experiences marginally greater declines in performance, it still delivers results that are largely consistent with expectations across the majority of tasks.

\subsection{Ablation Results}

\textbf{Effectiveness of the second positional interpolation}. As part of our progressive extension approach, we apply our search-based algorithm to perform a secondary non-uniform positional interpolation on LLMs that have already been fine-tuned and extended. To assess the benefits of this method, we conduct experiments utilizing our fine-tuned LLaMA2-256k model. We further expand this model to 512k, 1024k, and 2048k contexts, comparing our approach with PI and YaRN. The results in Table~\ref{tbl:secondsearch} indicate that our method maintains stable perplexity across various extension ratios. By contrast, both PI and YaRN experience a marked rise in perplexity as the extension ratio increases.

\textbf{Effectiveness of recovery at shorter context lengths.} In order to address the decline in model performance at reduced context lengths, we recalibrate the RoPE factors in LongRoPE-2048k using our search procedure. In this approach, we constrain the upper bound of the scaling factors during the search, which limits the degree of interpolation when working with shorter contexts such as 4k and 8k tokens. Table~\ref{tbl:shortperformance} presents a perplexity evaluation of LongRoPE-LLaMA2-2048k on Proof-pile for 4k and 8k contexts, alongside mean accuracy scores on LLM benchmarks. The data reveal that these adjustments lead to substantial improvements in model performance when handling shorter context lengths.

\textbf{Examination of the two types of non-uniformities}. To evaluate the individual impact of the two non-uniformities on model performance, we conduct ablation studies with two experimental designs: (i) for LLaMA2-7B, we extend context lengths to short 16k and 32k scenarios using various approaches—PI, optimising only the RoPE dimension, and optimising both forms of non-uniformity; (ii) for our already fine-tuned LLaMA2 model at 256k, we increase the context window to 2048k, applying the same strategies. We assess perplexity in both cases without any further fine-tuning. According to the results in Table~\ref{tbl:searchdimension}, adjusting for non-uniformity in the RoPE dimension offers a substantial reduction in perplexity compared with the PI method, which relies on linear interpolation. Additionally, introducing non-uniformity in token position yields noticeable gains at the 16k and 32k context lengths; however, this effect diminishes at the extensive 2048k length—likely due to the sheer scale involved. At this length, maintaining only the earliest tokens through non-interpolative means proves ineffective, highlighting a limitation to be addressed in future research.

\section{Related Works}

In addition to methods based on position interpolation, this section discusses related works of other approaches. 

\textbf{Retrieval-based} methods incorporate an external memory component to store extensive historical context and employ retrieval mechanisms to access pertinent documents during inference~\cite{longllama,wang2023augmenting,borgeaud2022improving}. Such methods usually require deliberate and significant alterations to the LLM architectures. By comparison, our approach offers a more streamlined solution, relying only on minimal adjustments to positional embeddings. Moreover, our methodology can be applied to a broader range of long-context tasks beyond retrieval, including long document summarization and few-shot learning.

\textbf{Attention-based context window extensions}. In addition to approaches utilizing positional embedding interpolation, several studies extend the input context window of the original LLM by altering attention mechanisms~\cite{han2023lminfinite, streamingllm,parallelwindow}. Central to this strategy is addressing the problem of attention explosion at novel positions through the introduction of innovative attention masks. Notably, these methods are compatible with and enhance the capabilities of positional interpolation techniques.

\textbf{Fine-tuning based approaches} aim to enhance the ability of pre-trained LLMs to handle extended contexts by adjusting position embeddings. For instance, models such as Code LLaMA~\cite{codellama}, LLaMA2 Long~\cite{xiong2023effective}, and ScaledRoPE~\cite{liu2023scaling} employ an enlarged base for RoPE and perform fine-tuning at the intended sequence lengths. In contrast, our approach supports flexible adaptation to a range of target lengths, reaching up to beyond 2M tokens. Given that extending context lengths beyond 128k via fine-tuning requires extensive GPU capabilities, recent methods like LongLoRA~\cite{longlora} and PoSE~\cite{zhu2023pose} have been designed to alleviate these resource demands. Our method complements these efficient fine-tuning strategies rather than overlapping with them.

\section{Conclusion}

This paper introduces LongRoPE, a technique that significantly increases the context window of LLMs to an unprecedented 2048k, all while preserving their proficiency within the original, shorter context limits. Our approach utilizes two distinct types of non-uniformities in RoPE positional embeddings, identified through an efficient evolutionary search process. This dual strategy confers two main advantages: it supplies robust initialization for subsequent fine-tuning, and it facilitates an eightfold expansion of the context window even in the absence of fine-tuning. Furthermore, we introduce a progressive extension method wherein models fine-tuned on a 256k-length context are used to extend up to a 2048k context window without necessitating additional fine-tuning. Comprehensive experimental results substantiate the efficacy of LongRoPE. We anticipate that our LongRoPE-2048k models will open up novel long-context use cases and catalyze further innovations in this domain.

\section*{Broader Impacts} This research aims to further the development of the Machine Learning discipline. While the study may have various implications for society, we do not identify any particular issues that require special emphasis at this time.

	\balance

\end{document}